%!TEX root = ../report.tex
\documentclass[report.tex]{subfiles}
\begin{document}

\begin{abstract}
Segmenting every frame of a UAV video with a per-frame model is expensive: the
strongest model trained here costs \LatTypical~ms per \SI{4}{K} frame, so a
\SI{45}{s} clip at \SI{20}{FPS} takes \ClipMinutes~minutes of GPU time.
Consecutive frames are almost identical, so most of that computation repeats
work already done.

This report evaluates an \emph{adaptive hybrid} alternative: segment one frame
with a model, propagate the mask to later frames by warping it along optical
flow, and re-run the model only when the propagated mask becomes unreliable.
Reliability is measured by a forward-backward consistency check, accumulated
across steps into a per-frame \emph{coverage} value; the model is re-invoked
when coverage falls to a threshold $\tau$, which resets coverage. Components
were selected by measurement: of the backbones trained behind one UPerNet
decoder, ConvNeXt-Large is strongest at \MIoUConvnextLarge{} mIoU, and SEA-RAFT
propagates as accurately as FlowNet2 at $3.6\times$ the speed.

On all \NumSeq{} UAVid validation sequences (\NumFrames{} frames), the pipeline
is \Speedup$\times$ faster than a model-every-frame baseline, retains
\Retention\% of its mIoU, and calls the model on \PctModel\% of frames under
one fixed threshold. Because flow is computed for every consecutive pair before
the threshold can be tested, flow is the fixed cost and model inference the
marginal one, so $\tau$ controls quality much more than speed. The residual
loss concentrates on small moving objects, where motion breaks the consistency
check.
\end{abstract}

\end{document}
